\documentclass[a4paper,12pt]{extreport}
% \usepackage{russianb}
\usepackage[russian]{babel}
\usepackage{float}
\usepackage[TS1,T2A]{fontenc}
\usepackage[utf8]{inputenc}

\usepackage{amsmath}
\usepackage{amssymb}
\usepackage[pdftex]{graphicx}
\graphicspath{{images/}}
\DeclareGraphicsExtensions{.eps,.pdf,.png,.jpg}

\usepackage{listings}
\usepackage{titlesec}

%Настройка оглавления
\titleformat{\chapter}[display]
    {\filcenter\large}
    {}
    {4pt}
    {\bfseries}{}

\titleformat{\section}
    {\normalsize\bfseries}
    {\thesection}
    {1em}{}

\titleformat{\subsection}
    {\normalsize\bfseries}
    {\thesubsection}
    {1em}{}

\titleformat{\subsubsection}
    {\normalsize\bfseries}
    {\thesubsubsection}
    {1em}{}

\titlespacing*{\chapter}{0pt}{30pt}{8pt}
\titlespacing*{\section}{\parindent}{*4}{*4}
\titlespacing*{\subsection}{\parindent}{*4}{*4}
\titlespacing*{\subsubsection}{\parindent}{*4}{*4}

%Задаем размер страницы
\usepackage{geometry}
\geometry{left=3cm}
\geometry{right=1cm}
\geometry{top=1cm}
\geometry{bottom=2cm}
\parindent=1cm

\makeatletter
\renewcommand\chapter{\par
\thispagestyle{plain}
\@afterindentfalse \secdef\@chapter\@schapter}
\makeatother

%Описание кода другим шрифтом.
\newcommand{\codepart}[1]
{{\ttfamily \textbf{#1}}}

\begin{document}
  \hfill
\begin{center}
  \hfill \break
  \includegraphics[scale=0.5]{dvfu_label}\\
  \normalsize{МИНИСТЕРСТВО ОБРАЗОВАНИЯ И НАУКИ}
  \normalsize{РОССИЙСКОЙ ФЕДЕРАЦИИ}\\
  \small{Федеральное государственное автономное образовательное учреждение}\\
  \small{высшего образования}\\
  \large{\textbf{«Дальневосточный федеральный университет»}}\\
  % \hfill \break
  \rule[0.0cm]{1\linewidth}{0.1cm}
  \rule[0.6cm]{1\linewidth}{0.05cm}
  % \hfill \break
  \large{\textbf{ШКОЛА ЕСТЕСТВЕННЫХ НАУК}}\\
  \hfill \break
  \large{\textbf{кафедра информатики, математического\\ и компьютерногомоделирования}}\\
  \hfill \break
  \hfill \break
  \large{ТЕМА КУРСОВОЙ РАБОТЫ}\\
  \hfill \break
  \hfill \break
  \hfill \break
  \large{\textbf{КУРСОВАЯ РАБОТА}}\\
  \large{по основной образовательной программе подготовки бакалавров\\
  по направлению 01.03.02 --- «Прикладная математика и информатика»\\
  профиль «Системное программирование»}\\
  % \hfill \break
  % \hfill \break
  % \hfill \break
  % \hfill \break
  \hfill \break
  \hfill \break
  \begin{tabular}{lcl}
    & & \small{Студент группы №\,БNNNNа}\\
    & & $\underset{\text{(подпись)}}{\underline{\hspace{4cm}}}$ \ Иванов И. И. \\
    & & «\underline{\hspace{1.25cm}}»\ \underline{\hspace{4cm}}2017г. \\
    Работа защищена & & \\
    с оценкой \hspace{1.3cm}\underline{\hspace{3.5cm}} & & Руководитель \\
    & & $\underset{\text{(должность, ученое звание)}}{\underline{\hspace{7cm}}}$ \\
    % & & \underline{\hspace{7cm}} \\
    & & $\underset{\text{(подпись)\hspace{2.5cm}(ФИО)}}{\underline{\hspace{7cm}}}$ \\
    Регистрационный номер\underline{\hspace{2cm}} & & «\underline{\hspace{1.25cm}}»\ \underline{\hspace{4cm}}2017г. \\
    %  & & \\
    «\underline{\hspace{1.25cm}}»\ \underline{\hspace{4cm}}2017г. & & \\
    & \hspace{0.5cm} & \\
  \end{tabular}
\end{center}

% \hfill \break
\hfill \break
\hfill \break
\hfill \break
\begin{center} \textbf{г. Владивосток\\ 2017} \end{center}
\thispagestyle{empty} % выключаем отображение номера для этой страницы

  \newpage
  \tableofcontents %Вывод оглавления
  \newpage
  % Обратите внимание на то, что неободимо задавать структуру с помощью chapter,
  % section, subsection и т.д. для корректной генерации оглавления
  \chapter{Основная часть}
\section{Пример Кода}
\noindent\indent Для того, что бы первый абзац в блоке был с красной строки необходимо написать \codepart{//noindent//indent}
\section{Секция с подсекциями}
\subsection{Подсекция 1}
\noindent\indent Для того, что бы первый абзац в блоке был с красной строки необходимо написать \codepart{//noindent//indent}
\subsection{Подсекция 2}
\noindent\indent Для того, что бы первый абзац в блоке был с красной строки необходимо написать \codepart{//noindent//indent}

\chapter{Еще одна глава}
\section{Пример Кода}
\noindent\indent Для того, что бы первый абзац в блоке был с красной строки необходимо написать \codepart{//noindent//indent}
\section{Секция с подсекциями}
\subsection{Подсекция 1}
\noindent\indent Для того, что бы первый абзац в блоке был с красной строки необходимо написать \codepart{//noindent//indent}
\subsection{Подсекция 2}
\noindent\indent Для того, что бы первый абзац в блоке был с красной строки необходимо написать \codepart{//noindent//indent}

\chapter*{Глава без номера, но указываемая в оглавлении}
\addcontentsline{toc}{chapter}{Глава без номера, но указываемая в оглавлении} %Использование внешнего файла с текстом

\end{document}
